\chapter{متغیر و حافظه}
\label{ch:variables}

\begin{goals}
\begin{itemize}
  \item حافظه‌ی رایانه چگونه چیزها را نگه می‌دارد
  \item تعریف متغیر با \code{:=} و تغییر مقدار آن با \code{=}
  \item تفاوت مقدار ثابت و مقدار تغییرپذیر (\fcode{ناپایا})
  \item قاعده‌های نام‌گذاری و انتخاب نام خوب
  \item چند الگوی پرکاربرد: شمارنده، جمع‌کردن، جابه‌جایی
\end{itemize}
\end{goals}

\section{جعبه‌هایی با برچسب}

در فصل پیش هر عددی را همان‌جا که لازم بود نوشتیم: \fcode{چاپ \lr{\salamcodefont ۲۰۲۶ - ۲۰۰۸}}.
اما برنامه‌های واقعی باید چیزهایی را \emph{به خاطر بسپارند}: نام کاربر، مبلغ
خرید، امتیاز بازی، تعداد دفعاتی که کاری انجام شده. جایی که برنامه این
چیزها را نگه می‌دارد، \textbf{حافظه}‌ی رایانه است.

حافظه را مثل قفسه‌ای بزرگ از جعبه‌های کوچک تصور کنید. در هر جعبه می‌توان
یک چیز گذاشت: یک عدد، یک متن، و یا چیزهایی که بعدها می‌بینیم. برای اینکه
بعداً بتوانیم جعبه را پیدا کنیم، روی آن برچسبی با یک نام می‌چسبانیم. به
چنین جعبه‌ای \textbf{متغیر} می‌گویند: جایی نام‌دار در حافظه که یک مقدار در
آن نگه داشته می‌شود.

\section{تعریف متغیر}

در سلام، برای ساختن یک جعبه‌ی تازه و گذاشتن چیزی در آن، از علامت \code{:=}
استفاده می‌کنیم. آن را «دونقطه مساوی» بخوانید، یا بهتر از آن: «تعریف کن
برابر با»:

\program{نخستین متغیرها}
\begin{salamfa}
روال ریشه:
    نام := "سارا"
    سن := ۱۷
    سرچاپ "نام:", نام
    سرچاپ "سن:", سن
پایان
\end{salamfa}

\begin{outputfa}
نام: سارا
سن: 17
\end{outputfa}

خط \fcode{نام \lr{:=} "سارا"} سه کار می‌کند: جعبه‌ای می‌سازد، برچسب
\fcode{نام} را روی آن می‌چسباند و متن \fcode{"سارا"} را در آن می‌گذارد. از
این پس، هر جا در برنامه \fcode{نام} بنویسیم، سلام به آن جعبه نگاه می‌کند و
مقدارش را برمی‌دارد. به همین دلیل \fcode{سرچاپ "نام:", نام} عبارت «نام:
سارا» را چاپ کرد، نه «نام: نام».

\begin{note}
به تفاوت \fcode{"نام"} و \fcode{نام} دقت کنید. اولی، با نقل قول، یک متن
است که همان‌طور که هست چاپ می‌شود. دومی، بدون نقل قول، نام یک متغیر است و
سلام به جای آن مقدار درون جعبه را می‌گذارد.
\end{note}

متغیرها را می‌توان در محاسبه به کار برد، درست مثل عددها:

\program{محاسبه با متغیرها}
\begin{salamfa}
روال ریشه:
    قیمت کتاب := ۱۲۰۰۰۰
    تعداد := ۳
    سرچاپ "مبلغ کل:", قیمت کتاب * تعداد, "تومان"
پایان
\end{salamfa}

\begin{outputfa}
مبلغ کل: 360000 تومان
\end{outputfa}

مزیت این کار روشن است: اگر فردا قیمت کتاب عوض شد، فقط یک خط را عوض می‌کنیم
و همه‌ی محاسبه‌هایی که از \fcode{قیمت کتاب} استفاده می‌کنند خودبه‌خود درست
می‌شوند.

\section{نام‌گذاری متغیرها}

در برنامه‌ی بالا نام یکی از متغیرها \fcode{قیمت کتاب} بود؛ دو واژه با یک
فاصله میانشان. این یکی از ویژگی‌های خوشایند سلام است: نام متغیر می‌تواند
از چند واژه ساخته شود، درست مثل زبان فارسی. قاعده‌های نام‌گذاری این‌هاست:
\begin{itemize}
  \item نام می‌تواند از حرف‌های فارسی، حرف‌های انگلیسی، رقم‌ها و فاصله
        ساخته شود.
  \item نام نمی‌تواند با رقم شروع شود: \fcode{۲نفر} نادرست است، اما
        \fcode{نفر۲} درست است.
  \item نام نمی‌تواند یکی از واژه‌های خود زبان باشد: نمی‌توانید متغیری به
        نام \fcode{سرچاپ} یا \fcode{پایان} یا \fcode{اگر} بسازید. این قاعده
        برای هر یک از واژه‌های یک نام چندواژه‌ای هم برقرار است؛ چون
        \fcode{هر}، \fcode{از}، \fcode{تا} و \fcode{و} واژه‌های
        زبان‌اند، نام‌هایی مثل \fcode{قیمت هر واحد} یا \fcode{بطری از جعبه}
        پذیرفته نمی‌شوند. فهرست این واژه‌ها را در پیوست \ref{app:keywords}
        می‌بینید.
  \item حرف بزرگ و کوچک در حرف‌های انگلیسی فرق دارد؛ \fcode{x} و
        \fcode{X} دو متغیر جداگانه‌اند.
\end{itemize}

مهم‌تر از قاعده‌ها، \emph{عادت‌های خوب} است. نام متغیر باید بگوید درونش چه
چیزی است. مقایسه کنید:

\begin{salamfa}
روال ریشه:
    الف := ۱۲۰۰۰۰
    ب := ۳
    سرچاپ الف * ب
پایان
\end{salamfa}

با همان برنامه‌ی پیشین. هر دو یک خروجی می‌دهند، اما دومی را یک هفته بعد هم
می‌توان فهمید و اولی را نه. برای نام‌های موقت و کم‌اهمیت (مثلاً در یک
محاسبه‌ی دو خطی) نام‌های کوتاه اشکالی ندارد، اما برای هر چیزی که در برنامه
می‌ماند، نام کامل و گویا انتخاب کنید.

\begin{tip}
در سلام می‌توانید نام متغیرها را انگلیسی هم بنویسید؛ مثلاً در برخی از
نمونه‌های ویرایشگر برخط، شمارنده‌ی حلقه \fcode{i} نام دارد. اما در این
کتاب همه‌ی نام‌ها را فارسی می‌نویسیم تا برنامه یک‌دست و روان خوانده شود.
\end{tip}

\section{تغییر دادن مقدار: مقدار ثابت و متغیر}

حالا نکته‌ای که سلام را از بسیاری زبان‌ها متمایز می‌کند. جعبه‌ای که با
\code{:=} می‌سازید، به طور پیش‌فرض \textbf{مهر و موم شده} است: مقداری که
اول در آن گذاشته‌اید، تا پایان همان می‌ماند. تلاش برای عوض کردن آن خطاست:

\begin{salamfa}
روال ریشه:
    سن := ۱۷
    سن = ۱۸
    سرچاپ سن
پایان
\end{salamfa}

\begin{outputfa}
خطا: نمی‌توان متغیر تغییرناپذیر 'سن' را دوباره مقداردهی کرد؛
برای این کار آن را 'ناپایا' تعریف کنید
 --> 3:5
\end{outputfa}

پیام خطا خودش راه حل را می‌گوید. اگر می‌خواهید مقدار جعبه‌ای در طول
برنامه عوض شود، هنگام تعریف، پیش از نامش واژه‌ی \kwd{ناپایا} را بنویسید:

\program{متغیر تغییرپذیر}
\begin{salamfa}
روال ریشه:
    ناپایا سن := ۱۷
    سرچاپ "امسال:", سن
    سن = ۱۸
    سرچاپ "سال بعد:", سن
پایان
\end{salamfa}

\begin{outputfa}
امسال: 17
سال بعد: 18
\end{outputfa}

به دو علامت متفاوت دقت کنید:
\begin{itemize}
  \item \code{:=} فقط \emph{یک بار}، هنگام ساختن جعبه، به کار می‌رود.
  \item \code{=} برای \emph{عوض کردن} مقدار جعبه‌ای است که از پیش وجود
        دارد و با \fcode{ناپایا} تعریف شده است.
\end{itemize}
اگر برای متغیری که پیش‌تر ساخته‌اید دوباره \code{:=} بنویسید، سلام خطا
می‌دهد که این نام از قبل تعریف شده است.

\begin{note}[چرا سلام این‌قدر سخت‌گیر است؟]
شاید بپرسید چرا باید زحمت نوشتن \fcode{ناپایا} را بکشیم. پاسخ این است که
بیشتر مقدارهای یک برنامه هرگز عوض نمی‌شوند: قیمت، نام، تعداد روزهای هفته.
وقتی این‌ها را بدون \fcode{ناپایا} تعریف می‌کنید، سلام تضمین می‌کند که هیچ
جای برنامه، حتی به اشتباه، آن‌ها را عوض نمی‌کند. و وقتی جایی \fcode{ناپایا}
می‌بینید، بلافاصله می‌فهمید «این یکی تغییر می‌کند، حواسم باشد». این عادت
کوچک، جلوی بسیاری از اشتباه‌های بزرگ را می‌گیرد.
\end{note}

\section{سلام هر متغیر را جدی می‌گیرد}

نکته‌ی دیگری که در همین ابتدا باید بدانید: اگر متغیری تعریف کنید و هیچ‌جا
از آن استفاده نکنید، سلام خطا می‌دهد:

\begin{salamfa}
روال ریشه:
    نام := "سارا"
    سن := ۱۷
    سرچاپ "نام:", نام
پایان
\end{salamfa}

\begin{outputfa}
خطا: متغیر استفاده‌نشده 'سن'
 --> 3:5
\end{outputfa}

دلیلش ساده است: متغیری که استفاده نمی‌شود یا اضافی است، یا نشانه‌ی این که
جایی چیزی را فراموش کرده‌اید. در هر دو حالت بهتر است بدانید. راه حل هم
ساده است: یا از آن استفاده کنید، یا خطش را پاک کنید.

\section{کپی شدن مقدارها}

وقتی مقدار یک متغیر را در متغیر دیگری می‌گذارید، \emph{یک نسخه} از آن مقدار
در جعبه‌ی دوم قرار می‌گیرد. دو جعبه به هم وصل نیستند:

\program{کپی شدن مقدار}
\begin{salamfa}
روال ریشه:
    ناپایا الف := ۱۰
    ب := الف
    الف = ۲۰
    سرچاپ "الف:", الف
    سرچاپ "ب:", ب
پایان
\end{salamfa}

\begin{outputfa}
الف: 20
ب: 10
\end{outputfa}

با اینکه \fcode{ب} از روی \fcode{الف} ساخته شد، عوض شدن \fcode{الف} هیچ
تأثیری روی \fcode{ب} نگذاشت. در لحظه‌ی \fcode{ب \lr{:=} الف}، عدد ۱۰ کپی
شد و رفت در جعبه‌ی \fcode{ب}؛ از آن به بعد هر جعبه راه خودش را می‌رود.

\section{الگوهای پرکاربرد}

چند کار هست که برنامه‌نویسان هر روز با متغیرها انجام می‌دهند. بیایید آن‌ها
را یکی‌یکی ببینیم.

\subsection{شمارنده}

شاید عجیب‌ترین خط برای تازه‌کارها این باشد:

\begin{salamfa}
شمارنده = شمارنده + ۱
\end{salamfa}

از نگاه ریاضی، این جمله بی‌معنی است؛ هیچ عددی با یک واحد بیشتر از خودش
برابر نیست. اما به یاد بیاورید که \code{=} در برنامه‌نویسی «برابر است» نیست،
بلکه «بگذار در» است. این خط می‌گوید: مقدار کنونی \fcode{شمارنده} را بردار،
یکی به آن اضافه کن و نتیجه را دوباره در همان جعبه بگذار. اول سمت راست
حساب می‌شود، بعد نتیجه در سمت چپ گذاشته می‌شود.

\program{شمارنده}
\begin{salamfa}
روال ریشه:
    ناپایا شمارنده := ۰
    سرچاپ شمارنده
    شمارنده = شمارنده + ۱
    سرچاپ شمارنده
    شمارنده = شمارنده + ۱
    سرچاپ شمارنده
پایان
\end{salamfa}

\begin{outputfa}
0
1
2
\end{outputfa}

این کار آن‌قدر رایج است که سلام برایش نوشتار کوتاه‌تری دارد.
\fcode{شمارنده \lr{+=} ۱} دقیقاً همان \fcode{شمارنده = شمارنده \lr{+} ۱}
است، و \fcode{شمارنده\lr{++}} همان \fcode{شمارنده \lr{+=} ۱}. هر سه شکل
درست‌اند و در این کتاب هر سه را خواهید دید:

\begin{salamfa}
روال ریشه:
    ناپایا شمارنده := ۰
    شمارنده = شمارنده + ۱
    شمارنده += ۱
    شمارنده++
    سرچاپ شمارنده
پایان
\end{salamfa}

\begin{outputfa}
3
\end{outputfa}

\subsection{جمع کردن}

همین الگو برای جمع زدن هم به کار می‌رود. متغیری با مقدار صفر می‌سازیم و هر
چیزی را که باید جمع شود به آن اضافه می‌کنیم:

\program{جمع خریدها}
\begin{salamfa}
روال ریشه:
    ناپایا مجموع := ۰
    مجموع += ۴۵۰۰۰    // نان
    مجموع += ۱۲۰۰۰۰   // پنیر
    مجموع += ۳۲۰۰۰    // شیر
    سرچاپ "جمع خرید:", مجموع, "تومان"
پایان
\end{salamfa}

\begin{outputfa}
جمع خرید: 197000 تومان
\end{outputfa}

\subsection{جابه‌جایی دو مقدار}

فرض کنید می‌خواهید محتوای دو جعبه را با هم عوض کنید. نمی‌شود اول \fcode{الف}
را در \fcode{ب} گذاشت و بعد \fcode{ب} را در \fcode{الف}، چون بعد از قدم
اول، مقدار قبلی \fcode{ب} از دست رفته است. راه حل، یک جعبه‌ی سوم است:

\program{جابه‌جایی}
\begin{salamfa}
روال ریشه:
    ناپایا الف := ۵
    ناپایا ب := ۹
    سرچاپ "پیش از جابه‌جایی:", الف, ب
    موقت := الف
    الف = ب
    ب = موقت
    سرچاپ "پس از جابه‌جایی:", الف, ب
پایان
\end{salamfa}

\begin{outputfa}
پیش از جابه‌جایی: 5 9
پس از جابه‌جایی: 9 5
\end{outputfa}

قدم به قدم دنبال کنید: \fcode{موقت} یک نسخه از ۵ را نگه می‌دارد، \fcode{الف}
مقدار ۹ می‌گیرد، و در آخر \fcode{ب} مقدار ۵ را از \fcode{موقت} پس می‌گیرد.
دقت کنید که \fcode{موقت} را بدون \fcode{ناپایا} تعریف کردیم، چون فقط یک بار
مقدار می‌گیرد و هرگز عوض نمی‌شود.

\section{پایدارها}

بعضی مقدارها نه تنها در طول اجرای برنامه عوض نمی‌شوند، بلکه اصلاً ذاتشان
ثابت است: تعداد روزهای هفته، نرخ مالیات، عدد پی. برای این‌ها سلام واژه‌ی
\kwd{پایا} را دارد:

\program{پایا}
\begin{salamfa}
روال ریشه:
    پایا نرخ := ۹
    قیمت := ۵۰۰۰۰۰
    مالیات := (قیمت * نرخ) / ۱۰۰
    سرچاپ "مالیات:", مالیات
    سرچاپ "قیمت با مالیات:", قیمت + مالیات
پایان
\end{salamfa}

\begin{outputfa}
مالیات: 45000
قیمت با مالیات: 545000
\end{outputfa}

از نگاه برنامه، \fcode{پایا نرخ \lr{:=} ۹} با \fcode{نرخ \lr{:=} ۹}
تفاوت چندانی ندارد؛ هر دو تغییرناپذیرند. تفاوت در پیام است: \fcode{پایا}
به خواننده می‌گوید «این یک عدد مهم است که با قصد این‌جا گذاشته‌ایم». دو
نکته‌ی کوچک: نام پایدار باید یک واژه باشد (برخلاف متغیرها که می‌توانند
چندواژه‌ای باشند)، و پایدارها را می‌توان بیرون از \fcode{روال ریشه} و در
بالای برنامه هم تعریف کرد؛ آن وقت در همه جای برنامه در دسترس‌اند. در فصل
\ref{ch:functions} این کار را خواهیم دید.

\begin{deep}[متغیر واقعاً کجاست؟]
حافظه‌ی رایانه واقعاً مثل قفسه‌ای از خانه‌های شماره‌دار است؛ به هر خانه یک
\textbf{نشانی} می‌گویند و در هر خانه یک عدد کوچک جا می‌شود. وقتی می‌نویسید
\fcode{سن \lr{:=} ۱۷}، مترجم سلام چند خانه‌ی خالی پیدا می‌کند، نشانی آن‌ها
را به نام \fcode{سن} گره می‌زند و عدد ۱۷ را در آن‌ها می‌نویسد. بعد از آن،
هر جا \fcode{سن} بنویسید، به آن نشانی ترجمه می‌شود. شما هرگز لازم نیست
نشانی‌ها را ببینید؛ نام‌ها برای همین اختراع شده‌اند. متن‌ها مثل
\fcode{"سارا"} کمی بیشتر جا می‌گیرند (هر حرف چند خانه)، اما اصل کار همان
است.
\end{deep}

\section{یک برنامه‌ی کامل‌تر}

صورت‌حساب یک خرید اینترنتی را با متغیرها بنویسیم:

\program{صورت‌حساب خرید}
\begin{salamfa}
روال ریشه:
    قیمت واحد := ۴۵۰۰۰
    تعداد := ۴
    هزینه ارسال := ۳۰۰۰۰
    جمع کالا := قیمت واحد * تعداد
    سرچاپ "جمع کالاها:", جمع کالا
    سرچاپ "هزینه ارسال:", هزینه ارسال
    سرچاپ "مبلغ قابل پرداخت:", جمع کالا + هزینه ارسال
پایان
\end{salamfa}

\begin{outputfa}
جمع کالاها: 180000
هزینه ارسال: 30000
مبلغ قابل پرداخت: 210000
\end{outputfa}

هیچ‌کدام از متغیرها با \fcode{ناپایا} تعریف نشده‌اند، چون هیچ‌کدام در طول
برنامه عوض نمی‌شوند. این برنامه را با نوشتن مستقیم عددها هم می‌شد نوشت،
اما آن وقت اگر تعداد از ۴ به ۵ تغییر می‌کرد، باید چند جا را دست می‌زدیم و
احتمالاً یکی را جا می‌انداختیم.

\begin{summary}
\begin{itemize}
  \item متغیر جایی نام‌دار در حافظه است که یک مقدار را نگه می‌دارد.
  \item \fcode{نام \lr{:=} مقدار} متغیر تازه‌ای می‌سازد که دیگر عوض نمی‌شود.
        \fcode{ناپایا نام \lr{:=} مقدار} متغیری می‌سازد که می‌توان بعداً با
        \code{=} مقدارش را عوض کرد.
  \item نام‌ها می‌توانند چندواژه‌ای باشند، اما با رقم شروع نمی‌شوند و
        واژه‌های زبان نیستند. نام گویا انتخاب کنید.
  \item متغیری که تعریف شده اما استفاده نشده، خطاست.
  \item \fcode{ب \lr{:=} الف} یک نسخه از مقدار می‌سازد؛ دو متغیر به هم وصل
        نیستند.
  \item \fcode{شمارنده \lr{+=} ۱} و \fcode{شمارنده\lr{++}} شکل‌های کوتاه
        \fcode{شمارنده = شمارنده \lr{+} ۱} هستند.
  \item \fcode{پایا} برای مقدارهایی است که ذاتاً ثابت‌اند.
\end{itemize}
\end{summary}

\section*{تمرین‌ها}

\exercise سه متغیر برای طول، عرض و ارتفاع یک اتاق (به متر) تعریف کنید و
مساحت کف و حجم اتاق را حساب و چاپ کنید.

\exercise متغیری به نام \fcode{موجودی} با مقدار اولیه‌ی ۵۰۰٬۰۰۰ بسازید.
سپس با سه دستور جداگانه، ۱۲۰٬۰۰۰ تومان برداشت، ۲۰۰٬۰۰۰ تومان واریز و
۷۵٬۰۰۰ تومان برداشت کنید و بعد از هر عملیات موجودی را چاپ کنید.

\exercise برنامه‌ی زیر را بدون اجرا کردن دنبال کنید و بگویید چه چاپ
می‌کند. سپس با اجرا کردن پاسخ خود را بررسی کنید:
\begin{salamfa}
روال ریشه:
    ناپایا الف := ۳
    ناپایا ب := الف * ۲
    الف = ب + ۱
    ب = الف - ب
    سرچاپ الف, ب
پایان
\end{salamfa}

\exercise برنامه‌ی زیر دو خطا دارد. آن‌ها را پیدا و برطرف کنید:
\begin{salamfa}
روال ریشه:
    امتیاز := ۱۰
    امتیاز = امتیاز + ۵
    نام بازیکن := "رضا"
    سرچاپ "امتیاز:", امتیاز
پایان
\end{salamfa}

\exercise سه متغیر \fcode{اول}، \fcode{دوم} و \fcode{سوم} با مقدارهای ۱، ۲
و ۳ بسازید و بدون نوشتن مستقیم عددها، مقدارها را طوری بچرخانید که
\fcode{اول} برابر ۲، \fcode{دوم} برابر ۳ و \fcode{سوم} برابر ۱ شود.
